MedGemma was the more accurate model (balanced accuracy 28.0\% on Derm7pt and 43.4\% on PAD-UFES-20, against 20.9\% and 27.7\% for Qwen2.5-VL; chance 12.5\%), and both models called about half of all lesions melanoma. With the lesion painted out, Qwen2.5-VL answered ``cannot be determined'' for 96.7\% of images, but MedGemma still gave a confident diagnosis (top probability at least 50\%) for 50.7\% (95\% CI 45.7--56.1), most often ``nevus'', and its mean confidence fell only from 75.7\% to 70.4\%. Changing only the skin around an unchanged lesion altered MedGemma's diagnosis in 6.2--13.4\% of images and Qwen2.5-VL's in 19.0--29.1\%; a ruler raised MedGemma's probability of melanoma by 1.9 points (1.1--2.8). MedGemma's yes/no answers on the 7-point checklist matched the expert labels better (AUROC 0.52--0.84) than Qwen2.5-VL's (0.46--0.74), and Qwen2.5-VL said ``yes'' to some criteria for almost every image. MedGemma was strongly overconfident (expected calibration error 0.35 and 0.39; mean confidence 68--81\% against accuracy 33--42\%), and so were the confidences both models stated in words.

\textbf{Conclusions.} The medical model is better at reading the lesion but does not know when there is no lesion to read and is far more confident than it is right; the general model is unstable under irrelevant changes. Neither behaviour shows up in an accuracy number. Counterfactual tests like these are cheap and should be reported alongside accuracy for medical VLMs.
